% !TeX root = ../papers/nc-p3.tex
\section{第三階数5で強制される低階数行列の超曲面}
\label{sec:rank-five-pencils}
本節では標準AS対象で$\rho=5$と仮定する。
$H$を五次元の係数空間、$A:H\to V_0\otimes V_2$と
$B:H^*\to V_1\otimes V_3$を最小分解による単射とする。
位置0の共変ベクトル$x$の評価$P(x):H\to V_2$の階数が2なら、
$W_x=(\ker P(x))^\perp\subset H^*$は二次元で、
$w(x)$の中央を縮約する行列束は$B(W_x)$である。他の位置も同様に読む。
その両端支持の次元がともに3以下となる階数2点の集合を$E_i$とする。
補題\ref{lem:middle-binary-pencil}より、このとき実際の束は転置を除き
$L_1$または$L_2$であり、両端の次元は$(1,2)$または$(2,3)$である。

\begin{lemma}[評価階数3の点は三文字列の中央に来ない]\label{lem:rho-five-rank-three-neighbors}
$r_i(y)=3$なら、各側の隣接零点は高々一つで、両側が同時に非空になることもない。
また各$(x,y)\in\Gamma_j$には$E_j$または$E_{j+1}$に属する端点がある。
\end{lemma}
\begin{proof}
例えば$y$を位置1とする。左の隣接点$x$は$r_0(x)=2$で、
$W_x$は位置1の評価の核（$H^*$内の二次元空間）に等しい。
命題\ref{prop:rank-two-evaluation-fiber}より$x$は一意である。
右の隣接点$z$についても同じで、その行空間$W_z$は同じ平面になる。
両方存在すると仮定すると$b_1(y)=f_1(y)=3$である。
$x,z$でそれぞれ縮約した三テンソルの位置1,3の支持は、いずれも
$B(W_x)=B(W_z)$のその両端の支持$L\subset V_1$, $N\subset V_3$である。
$y$が$L$を消すので$\dim L\le3$であり、
$b_0(x)+f_1(y)\le6$より$\dim N=b_0(x)\le3$である。
中央二次元の束の両支持次元は等しくない。そこで元の一次空間を順に
$k_x,V_1/L,k_z,V_3/N$へ射影すると全ての隣接関係が消え、
長さ$4m$の経路商次元は$((4-\dim L)(4-\dim N))^m\ge2^m$となる。矛盾である。

最後の主張を示す。$\Gamma_j$上で$f_j(x),b_{j+1}(y)\le3$である。
両方の評価階数が2で、どちらの端も$E$に入らなければ
$b_j(x)=f_{j+1}(y)=4$となり、その和は6以下という不等式に反する。
片方、例えば$y$が階数3ならば、$x$は階数2である。
$y$には前の隣接点があるため次の隣接点はなく、$f_{j+1}(y)=4$である。
従って$b_j(x)\le2$となって$x\in E_j$である。逆順も同じである。
\end{proof}

\begin{proposition}[一つの点から動く階数2束を得る]\label{prop:rho-five-propagate-small}
$x\in E_i$なら、その片側の隣接零点に$E$に属する点の正次元族がある。
さらに$x$の評価の三次元核に属する係数行列は全て階数2以下である。
\end{proposition}
\begin{proof}
$i=0$として順序の反転も許す。束が$L_2$型なら$b_0(x)=3$, $f_0(x)=2$とする。
隣接点$y$は射影直線をなす。全ての$y$について$f_1(y)\le6-b_0(x)=3$なので
次の隣接点も存在する。上の補題より$r_1(y)=2$であり、前後の支持がともに3以下なので
この直線全体が$E_1$に入る。

$L_1$型なら$f_0(x)=1$, $b_0(x)=2$とする。唯一の前方支持を$ku\subset V_1$と書くと、
$w(x)=u\otimes M$で$M\in V_2\otimes V_3$は階数2である。
従って$R_1$は$ku$と$M$の位置2の支持とのテンソル積を含む。
$R_1$の$(V_1/ku)\otimes V_2$への像は高々四次元である。
隣接点全体は$\mathbb P(u^\perp)=\mathbb P^2$であり、その像の基底を$y$で評価して得る列が
$V_2$を張らない点では次の隣接点も存在する。列が四本ならその行列式は
平面四次式であり、恒等的に零の場合を含め零点の次元は1以上である。
四本未満なら平面全体が該当する。これらの点は再び$E_1$に属する。

いずれも$y$の二次元の係数平面$W_y$は$K=\ker P(x)\subset H$に含まれる。
異なる$y$には異なる平面が対応する。各直線$\mathbb P(W_y)\subset\mathbb P(K)$上の
$A$行列は$L_1$または$L_2$型なので階数2以下である。
射影平面内の無限個の異なる直線は稠密であり、階数条件は閉であるため
$\mathbb P(K)$全体で階数2以下になる。
\end{proof}

\begin{theorem}[二つの係数空間の三次小行列式に共通因子が必要]\label{thm:rho-five-two-divisors}
標準AS対象で$\rho=5$なら、$A$と$B$の\emph{それぞれ}の射影四次元係数空間で、
階数2以下の部分は射影平面の動く族が掃く少なくとも三次元の閉部分を含む。
従ってどちらの側でも、全三次小行列式が恒等的に零であるか、
その非零なもの全てに共通の非定数既約因子がある。
特に各側の最大階数が4なら、その側で得た因子の次数は1または2である。
\end{theorem}
\begin{proof}
どの辺の$\Gamma$も非空で、少なくとも一端に$E$の点がある。
前命題をその点に適用すると他方の偶奇の位置に$E$の無限族が得られ、
そこから一点を選んで再び適用すると元の偶奇にも無限族が得られる。
位置は有限個であるから各偶奇で一位置を固定できる。
$E_i$は行列の階数条件で構成可能であり、正次元の既約な部分族を選べる。
各点の評価の核は命題\ref{prop:rank-two-evaluation-fiber}により互いに異なる。
階数2以下の閉部分が無限個の異なる射影平面を含むので、その次元は少なくとも3である。
実際、次元2以下の閉集合に含まれる平面はその既約成分に限られる。
射影四空間の既約な余次元1成分は一つの既約因子で定まり、
全ての三次小行列式はそれで割り切れる。
行列式が非零多項式である場合、その微分も三次小行列式の線型結合なので
同じ因子で割り切れる。標数零により四次の行列式はその因子の二乗で割り切れ、次数の上界を得る。
\end{proof}

\subsection{低階数因子の次数と正規形}
\begin{lemma}[階数2以下の係数空間と平面族が掃く超平面の排除]\label{lem:rank-five-hyperplane}
上の仮定の下では係数空間全体が階数2以下になることはない。
また前定理の$\mathbb P(\ker P(x))$の動く族の閉包が超平面になることもない。
\end{lemma}
\begin{proof}
まず一般行列の階数が2の場合を考え、係数行列の一つを$M=\operatorname{diag}(I_2,0)$とする。
$M+tA$のSchur補が恒等的に零なので、全ての係数行列は
$A=\left(\begin{smallmatrix}P&Q\\R&0\end{smallmatrix}\right)$であり、$RQ=0$となる。
$Q$が一般に可逆なら$R$は恒等的に零で一因子階数4に反する。
従って$Q$の線型空間は階数1以下である。階数1以下の線型空間は共通の像直線を持つか、
全元が一つの共変ベクトルとの外積である。実際、二つの純テンソルの和が純テンソルになるには
一方の因子が比例しなければならず、線型空間全体でも共通の側が決まる。
従ってある$0\ne v\in\ker M$で$Qv$の動く空間は一次元以下となり、$r_i(v)\ge2$に反する。
一般階数が1以下の場合も同じ純テンソルの事実が一因子階数4に反する。

次に動く平面の掃く既約超曲面が$\mathbb P(S)$、$S\subset H$が超平面であるとする。
対応する位置を0とする。動く点$x$で$S$の評価像は一次元である。
実際、$\ker P(x)\subset S$であり、全評価の階数は2である。既約な曲線族上で
その像直線を$ku_x$とする。一般の$z\in S$について$A(z)x$は一般の$x$で非零であり、
$u_x$の曲線の線型包は$\operatorname{im}A(z)$に入る。
その線型包が二次元以上ならば$S$上の階数上界2から、全$A(z)$の像が同じ平面に入る。
この平面を消す共変ベクトルで$H$を評価すると階数は高々1となり矛盾する。
従って$ku_x=ku$は固定である。点$x$の線型包を$U$とすれば$\dim U\ge2$で、
$S$は$U$を$ku$に送る。一般の$A(z)$は$U$上で非零なので、
$V_0^*/U\to V_2/ku$の誘導行列は全$S$上で階数1以下である。
共通像直線なら$A(S)$の像全体が一平面に入るため、上と同じ矛盾になる。
よって共通の共変ベクトルがあり、三次元の超平面$U'\supset U$が$S$により$ku$へ送られる。
誘導行列が全て零の場合も共通像平面の矛盾で除かれる。

$S=\ker\ell$とする。$U'$上では$z\in S$の作用を$A(z)v=f(v)(z)u$と書ける。
$f:U'\to S^*$には核がない（核点では全評価が階数1以下になる）。
$f$を$H^*$に線型に持ち上げれば、$0\ne v\in U'$について
$W_v=\langle\ell,f(v)\rangle$である。実際、$S$の補方向の作用が$ku$に入れば評価階数は1になる。
これらの直線は射影三空間$\mathbb P(\langle\ell,f(U')\rangle)$を覆う。
各$B(W_v)$は中央二次元の束なので、全点で$B(\ell)$と同じ階数である。
こうして$B$の超平面上に共通核・共通余核のない定数階数の$4\times4$行列空間を得る。
共通核や共通余核があれば残る一係数だけで評価することになり、階数下界2に反する。

この定数階数は1,2,3のいずれも不可能である。階数1は純テンソルの性質、
階数2は定理\ref{thm:rho-four-positive-incidence}の$\mathbb P^3$上の束の議論による。
階数3なら$4\mathcal O(-1)\to4\mathcal O$の核と余核は直線束となる。
それらの次数を$k,c$とするとChern指標から$c-k=4$, $c^2-k^2=-4$で、
$c+k=-1$となり整数解がない。これで超平面の場合を除いた。
\end{proof}

\begin{proposition}[係数の五次元空間の正規形]\label{prop:rho-five-spin-form}
$\rho=5$の標準AS性を仮定すると、各係数空間は左右および係数の基底変更で
\[
 A(S,n)=\begin{pmatrix}S&nI_2\\0&\operatorname{adj}S\end{pmatrix},
 \qquad S\in\operatorname{Mat}_{2\times2}(k),\quad n\in k
\]
に限られる。
\end{proposition}
\begin{proof}
上の平面族を含む既約超曲面の方程式を$F$とする。前補題から$\deg F\ge2$であり、
全随伴行列が零になる場合も除かれた。一般階数が3なら$\operatorname{adj}A=F L$で、
$L$は定数または一次の行列、一般に階数1である。
この線型空間（定数の場合はその張る直線）には固定像または固定共変ベクトルがある。
$AL=LA=0$より元の全係数にも共通の核または余核が生じ、一因子階数に反する。
従って一般階数は4で、前定理より$\det A=c F^2$、$F$は既約二次式である。
その階数は3または4である。実際、平面を含むので五変数非退化二次式ではなく、既約性から階数は3以上である。

$F(e)=1$とし、左基底変更で$A(e)=I$とする。$b$を$F$の極化で$b(e,e)=1$とする。
一般の$z$について$F(z-te)$は異なる二根を持ち、それぞれで$A(z)-tI$の核次元は2以上である。
従って二つの固有空間が全空間を張り、密度から全ての$z$で
\[
 A(z)^2-2b(e,z)A(z)+F(z)I=0 .
\]
$e^\perp$の非退化な直交方向の像は、適当にスカラーを掛けると二乗が恒等の反可換行列になる。
根基の像はそれらと反可換で二乗零である。

まず二つの非退化方向を取り、その像を$g_1,g_2$とする。
$g_2$が$g_1$の二固有空間を交換するので両方二次元であり、基底を選べば二行列は
$\operatorname{diag}(1,-1)\otimes I_2$と
$\left(\begin{smallmatrix}0&1\\1&0\end{smallmatrix}\right)\otimes I_2$である。
両者と反可換な行列は$g_1g_2(I_2\otimes T)$の形に限られる。
$F$が階数3なら根基は二次元であり、係数の単射性から二次元の$T$の空間の全元が二乗零になる。
しかし二行二列の跡零行列の射影平面で行列式零は滑らかな二次曲線なので直線を含まない。

従って$F$は階数4である。第三の非退化方向$g_3$も上の$T$を使って表され、
$g_3^2=I$より$T$の二乗は非零スカラーである。
根基の非零な像はこれと反可換であるため、この$T$はスカラーではない。
その二固有空間に従って四次元を二つの二次元ブロックに分けると、
初めの二行列の作用は両ブロックで同じ、第三の作用は符号が逆である。
第二ブロックを$g_1g_2$の二次元因子で共役すれば三行列全てが第一ブロックの負になる。
三行列は第一ブロックの跡零空間を張る。従って$ke$と三方向の作用は
$\operatorname{diag}(S,\operatorname{adj}S)$である。
根基は三つ全てと反可換なので対角ブロックは零、非対角ブロックはスカラーとなる
（最初の二行列との反可換行列はそれらの積の倍であり、第三とは可換する）。
二乗零と非零性より片側の非対角ブロックだけが非零であり、
ブロックの交換と根基の基底変更で表示を得る。
\end{proof}

\begin{theorem}[第三階数の鋭い下界]\label{thm:as-opposite-rank-sharp-six}
標準14641型の全AS対象の第三階数は$6\le\rho\le16$を満たす。
下界6は可換多項式環で達成される。
\end{theorem}
\begin{proof}
階数5を仮定して位置0,2の係数を前命題の表示にする。
$S=\left(\begin{smallmatrix}a&b\\c&d\end{smallmatrix}\right)$と記す。
位置0で入力$(v,w)$を固定して評価すると、$w=0$, $v=(s,t)\ne0$のとき階数2で、
その行空間は$\langle sa+tb,sc+td\rangle$である。
$w\ne0$なら$\operatorname{adj}S\,w$が二次元を張り、$n w$が上ブロックにもう一方向を与えるので
階数は3以上である。転置側は上ブロックの共変ベクトルが非零なら同じく階数3以上である。
下ブロックだけの共変ベクトル$(u,v')$では行空間は
$\langle ud-v'c,-ub+v'a\rangle$となる。
従って両側の階数2点は同じ二次元行空間の族を与え、$(s,t)=(v',-u)$で対応する。

命題\ref{prop:rho-five-propagate-small}と定理\ref{thm:rho-five-two-divisors}の証明により、
位置0または2には$E$の点がある。その反対側にも同じ行空間$W$の点がある。
$B(W)$の両端の支持の次元は$(1,2)$または$(2,3)$（順序の交換を許す）。
この二つの支持を除く商を位置1,3に取り、位置0,2を対応する二共変ベクトルで一次元に落とせば、
補題\ref{lem:rho-five-rank-three-neighbors}と同じく四辺の関係は全て零である。
長さ$4m$の商が少なくとも$2^m$次元となり、AS対象の多項式成長に反する。
階数4以下の排除は定理\ref{thm:as-opposite-rank-five-lower}、可換例の階数は
交代四テンソルの二対二切断の階数$\dim\bigwedge^2 k^4=6$である。
\end{proof}
